%-------------------------------------------------------------------------------
%	ABSTRACT PAGE
%-------------------------------------------------------------------------------

\addchaptertocentry{Хураангуй} % Хураангуйг гарчигт нэмэх

\begin{center}
{\scshape\Large \univname\par} % Их сургуулийн нэр
{\scshape\large \facname\par}\vspace{0.5cm} % Их сургуулийн нэр
{\huge\textbf{{Хураангуй}} \par}
\bigskip
{\Large{\ttitle} \par} % Тезисийн нэр
\bigskip

{\normalsize \shortname \par} % Зохиогчийн нэр
\addressname
\end{center}

\textit{\textbf{Түлхүүр үгс: \keywordnames}}
\bigskip

Өнөөгийн дижитал хөгжмийн үйлдвэрлэлд FL Studio зэрэг DAW програмууд өргөн ашиглагдаж байгаа боловч эхлэн суралцагчдын хувьд beat хийх, melody бичих, arrangement хийх, mixing/mastering хийх дараалал ойлгомжгүй хэвээр байна. Интернэт орчинд сургалтын материал олон байгаа ч тэдгээр нь нэгтгэсэн бүтэц, явцын хяналт, дадлага, мэргэжлийн зөвлөгөөний орчноор дутмаг байдаг.

Энэхүү төгсөлтийн ажлын хүрээнд FL Studio болон music production сурах хүсэлтэй хэрэглэгчдэд зориулсан төлбөртэй онлайн сургалтын вебсайтыг боловсруулж, хөгжүүлэв. Уг систем нь мэргэжлийн продюсер багш нарын бэлтгэсэн tutorial video, course catalog, paid lesson access, lesson player, dashboard, self-check quiz, progress tracking зэрэг сургалтын үндсэн функцуудыг агуулна. Сургалтын агуулга нь rhythm, melody, harmony, arrangement, beat making, mixing, mastering зэрэг чадварыг шат дараатай эзэмшүүлэх workflow болон дадлага дээр суурилна.

Хиймэл оюун ухааны хэсэг нь системийн үндсэн бүтээгдэхүүн бус, харин суралцагчийн асуултад туслах нэмэлт боломж юм. Тухайлбал, FL Studio, beat making, mixing, mastering болон хөгжмийн онолын асуултад хариулах RAG технологид суурилсан chatbot нь хичээл үзэх явцад ойлгоогүй зүйлээ лавлахад ашиглагдана. Иймээс энэхүү ажил нь мэргэжлийн продюсеруудын төлбөртэй видео хичээлийг бүтэцтэй хүргэх, суралцагчийн явцыг хянах, хэрэгцээтэй үед AI туслах зөвлөгөө авах боломжийг нэгтгэсэн сургалтын платформ болж байгаагаараа практик ач холбогдолтой.
